\documentclass[twocolumn,prl,superscriptaddress,nofootinbib]{revtex4-2}

\usepackage{amsmath,amssymb}
\usepackage{mathtools}
\usepackage{booktabs}
\usepackage{hyperref}

\newcommand{\R}{\mathbb{R}}
\newcommand{\C}{\mathbb{C}}
\newcommand{\Z}{\mathbb{Z}}
\newcommand{\Q}{\mathbb{Q}}
\newcommand{\KS}{\mathrm{KS}}

\begin{document}

\title{Kochen--Specker uncolorability in cyclotomic fields\\requires exactly $6 \mid n$}

\author{Michael Kernaghan}
\affiliation{Pacific Quantum Systems, Vancouver, Canada}

\date{\today}

\begin{abstract}
We prove that the ray set generated by the $n$th roots of unity in $\C^3$ is Kochen--Specker uncolorable if and only if $6 \mid n$.  Sufficiency follows from the embedding of Cabello's Eisenstein 33-vector KS set.  Necessity is proved by algebraic analysis of vanishing sums among roots of unity: when $n$ is odd and $3 \mid n$, a projective collapse argument shows every triad is isolated (no two share a ray), so the coloring problem decomposes; when $3 \nmid n$ and $2 \mid n$, the orthogonality graph on each coordinate plane is a perfect matching, admitting an explicit coloring.  The factorization $6 = 2 \times 3$ mirrors the result of Cortez, Morales, and Reyes that $\Z[1/6]$ is the minimal ring supporting KS contextuality in dimension~3: the arithmetic conditions $2 \mid n$ (two-term cancellation) and $3 \mid n$ (three-term phase cancellation) are independently necessary and jointly sufficient.
\end{abstract}

\maketitle

\section{Introduction}

The Kochen--Specker (KS) theorem~\cite{KochenSpecker1967} establishes that quantum measurements in dimension $d \geq 3$ cannot be explained by noncontextual hidden-variable theories, by exhibiting a finite set of rays admitting no consistent $\{0,1\}$-valued coloring.  A natural source of KS constructions is the cyclotomic family: for each positive integer~$n$, the alphabet $\{0\} \cup \{\zeta^k : k = 0, \ldots, n-1\}$ with $\zeta = e^{2\pi i/n}$ generates a ray set $S_n$ in $\C^3$ via all projectively distinct nonzero vectors with coordinates drawn from this alphabet.

Cabello~\cite{Cabello2025simplest} showed that $S_6$ contains a 33-vector KS set built from the Eisenstein integers $\Z[\omega]$ ($\omega = e^{2\pi i/3}$), which he identified as the ``simplest'' KS set by the criterion of fewest bases.  Cortez, Morales, and Reyes~\cite{CortezMoralesReyes2022} proved that $\Z[1/6]$ is the minimal ring extension of~$\Z$ exhibiting KS contextuality in dimension~3---that is, the primes 2 and~3 are both necessary.

These results suggest a sharp characterization: KS-uncolorability in the cyclotomic family should require exactly $6 \mid n$.  We prove this.

\begin{theorem}\label{thm:main}
The ray set $S_n$ generated by the $n$th roots of unity in $\C^3$ is KS-un\-color\-able if and only if $6 \mid n$.
\end{theorem}

\section{Preliminaries}

A \emph{ray} in $\C^3$ is a one-dimensional subspace.  Two rays $v, w$ are \emph{orthogonal} if $\langle v | w \rangle = \sum_k \bar{v}_k w_k = 0$.  An orthogonal \emph{triad} is a set of three mutually orthogonal rays.  A \emph{KS coloring} assigns each ray a value in $\{0,1\}$ such that orthogonal rays cannot both receive~1 and every triad contains exactly one~1.  A ray set is \emph{KS-uncolorable} if no such coloring exists.

For $\zeta = e^{2\pi i/n}$, the ray set $S_n$ consists of all projectively distinct nonzero vectors $(a_1, a_2, a_3)$ with each $a_k \in \{0\} \cup \{\zeta^j : 0 \leq j < n\}$.  The orthogonality condition $\langle v | w \rangle = 0$ reduces to a vanishing sum of terms, each of the form $\zeta^m$ or~$0$.

\section{Proof of Theorem~\ref{thm:main}}

\subsection{Sufficiency}

If $6 \mid n$, then $\zeta^{n/6}$ generates the 6th roots of unity as a subgroup of $\langle\zeta\rangle$, so the alphabet for $n = 6$ embeds into that for general~$n$.  Since $S_6$ contains Cabello's 33-vector KS set~\cite{Cabello2025simplest}, $S_n$ is KS-uncolorable. \qed

\subsection{Necessity: vanishing-sum analysis}

We show that if $6 \nmid n$, then $S_n$ admits a valid KS coloring.  The proof rests on two lemmas characterizing when vanishing sums of roots of unity exist.

\begin{lemma}[Three-term vanishing sums]\label{lem:3term}
The equation $\zeta^a + \zeta^b + \zeta^c = 0$ has a solution with $\zeta = e^{2\pi i/n}$ if and only if $3 \mid n$.
\end{lemma}

\begin{proof}
Dividing by $\zeta^a$: $1 + \zeta^p + \zeta^q = 0$ where $p = b - a$, $q = c - a$.  Setting $\alpha = \zeta^p$, we need $\zeta^q = -1 - \alpha$ with $|\zeta^q| = 1$.  Computing $|-1 - \alpha|^2 = 2 + 2\cos(2\pi p/n)$, the condition $|-1-\alpha| = 1$ forces $\cos(2\pi p/n) = -1/2$, hence $p/n \in \{1/3, 2/3\}$, which requires $3 \mid n$.  Conversely, if $3 \mid n$, then $1 + \zeta^{n/3} + \zeta^{2n/3} = 0$.
\end{proof}

\begin{lemma}[Two-term vanishing sums]\label{lem:2term}
The equation $\zeta^a + \zeta^b = 0$ has a solution if and only if $2 \mid n$.
\end{lemma}

\begin{proof}
$\zeta^a + \zeta^b = 0$ iff $\zeta^{b-a} = -1$ iff $-1 \in \langle\zeta\rangle$ iff $2 \mid n$.
\end{proof}

These lemmas constrain which orthogonalities can occur in $S_n$.  Two all-nonzero rays $v = (\zeta^{a_1}, \zeta^{a_2}, \zeta^{a_3})$ and $w = (\zeta^{b_1}, \zeta^{b_2}, \zeta^{b_3})$ are orthogonal iff $\zeta^{b_1 - a_1} + \zeta^{b_2 - a_2} + \zeta^{b_3 - a_3} = 0$---a three-term vanishing sum.  By Lemma~\ref{lem:3term}, this requires $3 \mid n$.  Similarly, orthogonality between one-zero rays in the same coordinate plane requires a two-term vanishing sum, hence $2 \mid n$ by Lemma~\ref{lem:2term}.

We now treat the three cases where $6 \nmid n$.

\subsection{Case 1: $3 \nmid n$, $2 \nmid n$}

No vanishing sums of any length exist among the $n$th roots.  The only orthogonalities are axis-complementary (e.g., $(\zeta^a, 0, 0) \perp (0, \zeta^b, \zeta^c)$).  The sole triad is $\{[1{:}0{:}0],\, [0{:}1{:}0],\, [0{:}0{:}1]\}$, which is trivially colorable. \qed

\subsection{Case 2: $3 \nmid n$, $2 \mid n$}

By Lemma~\ref{lem:3term}, no all-nonzero orthogonal pairs exist.  All triads have the form $\{v_1, v_2, e_k\}$ where $v_1, v_2$ are one-zero rays in the same coordinate plane and $e_k$ is the corresponding axis ray.

\begin{lemma}[Perfect matching]\label{lem:matching}
For even~$n$, the orthogonality graph on one-zero rays in each coordinate plane is a perfect matching.
\end{lemma}

\begin{proof}
The one-zero rays in the $z = 0$ plane are projectively $(1, \zeta^p, 0)$ for $p \in \{0, \ldots, n-1\}$.  Two such rays $p, q$ are orthogonal iff $1 + \zeta^{q-p} = 0$, i.e., $q - p \equiv n/2 \pmod{n}$.  Each ray has exactly one partner, giving $n/2$ disjoint edges.
\end{proof}

\begin{lemma}[Plane independence]\label{lem:planes}
One-zero rays from different coordinate planes are never orthogonal.
\end{lemma}

\begin{proof}
$\langle(\zeta^a, \zeta^b, 0) | (\zeta^c, 0, \zeta^d)\rangle = \zeta^{c-a} \neq 0$.
\end{proof}

An explicit coloring: choose one axis ray as green (say $e_1$), the other two as red.  The $e_1$-plane rays are all forced red (orthogonality exclusion).  In each of the other two planes, the matched pairs can be independently 2-colored (one green, one red per pair).  By Lemma~\ref{lem:planes}, no cross-plane conflicts arise. \qed

\subsection{Case 3: $3 \mid n$, $2 \nmid n$}

This is the principal case.  Three-term cancellations create triads among all-nonzero rays, but we show these triads are completely isolated.

\begin{lemma}[Projective collapse]\label{lem:collapse}
For any all-nonzero ray $v = (\zeta^a, \zeta^b, \zeta^c)$ in $S_n$ with $3 \mid n$, there are exactly two projectively distinct all-nonzero rays orthogonal to~$v$.
\end{lemma}

\begin{proof}
Let $\omega = \zeta^{n/3}$.  The orthogonality condition $\langle v | w \rangle = 0$ with $w$ all-nonzero requires the coordinate differences $(b_1 - a_1, b_2 - a_2, b_3 - a_3) \bmod n$ to be a permutation of $(0, n/3, 2n/3)$ (by Lemma~\ref{lem:3term}).  The six permutations split into two projective equivalence classes.

\emph{Even permutations} $\{(0, n/3, 2n/3),\, (n/3, 2n/3, 0),\, (2n/3, 0, n/3)\}$ yield:
\begin{align*}
w_+ &= (\zeta^a,\, \omega\zeta^b,\, \omega^2\zeta^c), \\
&= \omega^{-1} \cdot (\omega\zeta^a,\, \omega^2\zeta^b,\, \zeta^c), \\
&= \omega^{-2} \cdot (\omega^2\zeta^a,\, \zeta^b,\, \omega\zeta^c).
\end{align*}
All three are projectively equivalent (they differ by the scalar factors $1, \omega^{-1}, \omega^{-2}$).

\emph{Odd permutations} $\{(0, 2n/3, n/3),\, (n/3, 0, 2n/3),\, (2n/3, n/3, 0)\}$ similarly collapse to
\[
w_- = (\zeta^a,\, \omega^2\zeta^b,\, \omega\zeta^c).
\]
Since $\omega \neq 1$, neither $w_+$ nor $w_-$ is projectively equivalent to~$v$, and $w_+ \not\sim w_-$ (else $\omega = \omega^2$, forcing $\omega = 1$).
\end{proof}

\begin{lemma}[Triad uniqueness]\label{lem:unique}
The triple $\{v, w_+, w_-\}$ is a triad, and it is the unique triad containing~$v$ among all-nonzero rays.
\end{lemma}

\begin{proof}
$\langle w_+ | w_- \rangle = 1 + \bar\omega \cdot \omega^2 + \bar\omega^2 \cdot \omega = 1 + \omega + \omega^2 = 0$ (using $\bar\omega = \omega^2$).  So $\{v, w_+, w_-\}$ is mutually orthogonal.  Uniqueness follows from Lemma~\ref{lem:collapse}: $v$ has exactly two all-nonzero orthogonal neighbors.
\end{proof}

\begin{lemma}[One-zero isolation]\label{lem:onezero}
For odd~$n$, no one-zero ray participates in any triad.
\end{lemma}

\begin{proof}
A triad containing $({\zeta^a}, {\zeta^b}, 0)$ would require a second ray~$u$ with $u_3 = 0$ (forced by orthogonality to the complementary axis ray $(0, 0, \zeta^c)$) satisfying $\langle v | u \rangle = \zeta^{e-a} + \zeta^{f-b} = 0$, i.e., $-1 \in \langle\zeta\rangle$.  By Lemma~\ref{lem:2term}, this requires $2 \mid n$.
\end{proof}

The triad structure of $S_n$ for odd $n$ with $3 \mid n$ is now determined: one axis triad $\{e_1, e_2, e_3\}$ plus $n^2/3$ all-nonzero triads (since the $n^2$ all-nonzero rays partition into triples by Lemma~\ref{lem:unique}), with every triad isolated---no two sharing a ray.  The coloring problem decomposes into $n^2/3 + 1$ independent single-triad subproblems, each trivially satisfiable by assigning one~1 per triad. \qed

\section{Discussion}

The factorization $6 = 2 \times 3$ in Theorem~\ref{thm:main} reflects two independent algebraic mechanisms, each necessary for KS-uncolorability:

\emph{Phase cancellation} ($3 \mid n$): the identity $1 + \omega + \omega^2 = 0$ creates orthogonal triads among all-nonzero rays.  Without it, no all-nonzero pair is orthogonal, and triads are restricted to zero-coordinate configurations.

\emph{Interlocking} ($2 \mid n$): the element $-1 \in \langle\zeta\rangle$ creates two-term cancellations that link triads together.  Without it, the projective collapse of Lemma~\ref{lem:collapse} ensures every triad is isolated, making the coloring problem decomposable.

This mirrors the ring-theoretic result of Cortez, Morales, and Reyes~\cite{CortezMoralesReyes2022}: $\Z[1/6]$ is the minimal ring supporting KS contextuality in dimension~3 precisely because both $1/2$ and $1/3$ are needed in the denominators of the projection matrices $P_v = vv^\dagger / \|v\|^2$.  Computing their $N(S) = \operatorname{lcm}\{\|v\|^2 : v \in S\}$ invariant for the Eisenstein 33-set gives $N = 6$, confirming it realizes the minimal ring.  Theorem~\ref{thm:main} shows that this minimality is not an artifact of a particular construction but a structural feature of the cyclotomic family: among all $n$th-root alphabets, exactly those with $6 \mid n$ inherit both mechanisms.

More broadly, the proof technique---classifying orthogonalities by the length of the underlying vanishing sum, then analyzing the resulting constraint graph---may extend to other algebraic families of KS constructions.  The projective collapse phenomenon (Lemma~\ref{lem:collapse}), where group-theoretic symmetry forces orthogonal neighbors to coalesce, suggests a representation-theoretic approach to KS colorability that goes beyond the combinatorial methods currently dominant in the field.

\begin{acknowledgments}
The author acknowledges the use of Claude (Anthropic) for computational assistance and proof verification.  SAT solving used PySAT with Glucose4.  Code is available at \url{https://github.com/michaelkernaghan/contextuality}.
\end{acknowledgments}

\begin{thebibliography}{10}

\bibitem{KochenSpecker1967}
S.~Kochen and E.~P.~Specker,
J.\ Math.\ Mech.\ \textbf{17}, 59 (1967).

\bibitem{Cabello2025simplest}
A.~Cabello,
Phys.\ Rev.\ Lett.\ (2025); arXiv:2508.07335.

\bibitem{CortezMoralesReyes2022}
V.~Cortez, R.~Morales, and E.~J.~Reyes,
arXiv:2211.13216 (2022).

\bibitem{ConwayKochen}
J.~H.~Conway and S.~Kochen,
reported in A.~Peres, \textit{Quantum Theory: Concepts and Methods} (Kluwer, 1993).

\bibitem{TrandafirCabello2025}
S.~Trandafir and A.~Cabello,
Phys.\ Rev.\ A \textbf{111}, 052204 (2025); arXiv:2501.11640.

\bibitem{LiBrightGanesh2024}
Z.~Li, C.~Bright, and V.~Ganesh,
in \textit{Proceedings of IJCAI-24}, 2105 (2024); arXiv:2306.13319.

\end{thebibliography}

\end{document}
